%% ECON 282E -- Session 7, Deck B2: Deep VFI and Actor-Critic
%% Sources: AI-ECON Lec04_T2_VFI_DL_ActorCritic.tex, carried WHOLE (20 frames);
%% AI-ECON Lec06_T2a_Why_ML_Euler.tex fr 3-10; Quant_Macro Lec.6
%% "DNN Based Value Function Iteration" and "The Role & Advantage of
%% Reinforcement Learning".
%% Numbers from tools/figures/s07_numbers.json (key "five_methods") and
%% s04_numbers.json (key "optimizer_memory"), both cited, not recomputed.

%% ECON 282E -- Foundations of Macroeconomics (UC Riverside, Fall 2026)
%% Shared Beamer preamble for every deck in slides/S01 ... slides/S10.
%%
%% Usage (a deck lives in slides/S0X/ and is compiled from that folder):
%%   %% ECON 282E -- Foundations of Macroeconomics (UC Riverside, Fall 2026)
%% Shared Beamer preamble for every deck in slides/S01 ... slides/S10.
%%
%% Usage (a deck lives in slides/S0X/ and is compiled from that folder):
%%   %% ECON 282E -- Foundations of Macroeconomics (UC Riverside, Fall 2026)
%% Shared Beamer preamble for every deck in slides/S01 ... slides/S10.
%%
%% Usage (a deck lives in slides/S0X/ and is compiled from that folder):
%%   \input{../preamble/course_preamble.tex}
%%   \decktitle{1}{A1}{Who I Am \& Why This Course}{September 24, 2026}
%%   \begin{document}
%%   \makedecktitle
%%   ...
%%
%% Adapted from Courses/AI-ECON-2026/source/shared_preamble.tex (Metropolis theme,
%% concept/caution/example/takeaway boxes, \sectiondivider). Changes for this course:
%%   * course title block (\decktitle / \makedecktitle) instead of \makebeamertitle
%%   * \zh{...} typesets Chinese (book titles) under pdflatex via CJKutf8
%%   * researchbox for the "Research in the wild" frames
%%   * section pages off; TikZ libraries and the course map loaded here
%% Build with pdflatex only (two passes); tools/build_slides.py does this for you.

\documentclass[10pt,english,aspectratio=169]{beamer}

\usepackage[T1]{fontenc}
\usepackage{textcomp}
\usepackage[utf8]{inputenc}
\usepackage{babel}
\usepackage{CJKutf8}

\usepackage{amsmath, amssymb, amsthm, graphicx}
\usepackage{booktabs, multirow, tabularx}
\usepackage{tikz}
\usetikzlibrary{positioning,arrows.meta,shapes,calc,fit,backgrounds}
\usepackage{pgfplots}
\pgfplotsset{compat=1.16}
\usepackage[most]{tcolorbox}
\tcbuselibrary{skins, breakable}
\usepackage{listings}
\usepackage{xcolor}

\lstset{
  basicstyle=\ttfamily\small,
  keywordstyle=\color{blue!80!black}\bfseries,
  commentstyle=\color{green!50!black}\itshape,
  stringstyle=\color{orange!80!black},
  backgroundcolor=\color{gray!8},
  frame=single,
  rulecolor=\color{gray!40},
  breaklines=true,
  showstringspaces=false,
  numbers=none,
}

\usetheme{metropolis}
\usefonttheme{professionalfonts}
\metroset{sectionpage=none}

\definecolor{LightGray}{RGB}{230,230,230}
\definecolor{RedTitle}{RGB}{0,0,200}
\definecolor{VeryLightGray}{RGB}{245,245,245}
\definecolor{DarkText}{RGB}{0,0,0}
\definecolor{UCRBlue}{RGB}{0,45,106}
\definecolor{UCRGold}{RGB}{241,171,0}

% Stage colours of the course arc (used by the course map and elsewhere)
\colorlet{StageTrunk}{blue!12}
\colorlet{StageBranch}{cyan!14}
\colorlet{StageMachine}{gray!18}
\colorlet{StageWall}{red!14}
\colorlet{StageTools}{orange!20}
\colorlet{StageData}{green!16}
\colorlet{StageAgents}{violet!14}

\hypersetup{bookmarks=false,breaklinks=true,pdfborder={0 0 0},colorlinks=true,
  linkcolor=DarkText,urlcolor=RedTitle}

\setbeamercolor{background canvas}{bg=white}
\setbeamercolor{title}{fg=RedTitle}
\setbeamercolor{frametitle}{bg=VeryLightGray,fg=DarkText}
\setbeamercolor{normal text}{fg=DarkText}
\setbeamercolor{alerted text}{fg=RedTitle}
\setbeamersize{text margin left=5mm, text margin right=5mm}
\addtobeamertemplate{frametitle}{}{\vspace{-0.5em}}

\setbeamertemplate{navigation symbols}{}
\setbeamertemplate{footline}{%
  \hspace{5mm}{\usebeamerfont{page number in head/foot}\color{black!45}%
  ECON 282E \,$\cdot$\, \insertshortdeck}\hfill%
  \usebeamercolor[fg]{page number in head/foot}%
  \usebeamerfont{page number in head/foot}%
  \insertframenumber\,/\,\inserttotalframenumber\kern1em\vskip2pt%
}

% ---------------------------------------------------------------------------
% Chinese text under pdflatex (book titles etc.)
\newcommand{\zh}[1]{\begin{CJK*}{UTF8}{gbsn}#1\end{CJK*}}

% ---------------------------------------------------------------------------
% Deck title block
%   \decktitle{session}{deck id}{title}{date}
\newcommand{\insertshortdeck}{}
\newcommand{\decksession}{}
\newcommand{\deckid}{}
\newcommand{\decktitletext}{}
\newcommand{\deckdate}{}
\newcommand{\decktitle}[4]{%
  \renewcommand{\decksession}{#1}%
  \renewcommand{\deckid}{#2}%
  \renewcommand{\decktitletext}{#3}%
  \renewcommand{\deckdate}{#4}%
  \renewcommand{\insertshortdeck}{Session #1 \,$\cdot$\, Deck #1.#2}%
  \title{#3}%
  \author{Zhigang Feng}%
  \date{#4}%
}
\newcommand{\makedecktitle}{%
  \begin{frame}[plain,noframenumbering]
    \begin{tikzpicture}[remember picture,overlay]
      \fill[UCRBlue] (current page.north west) rectangle ([yshift=-0.55cm]current page.north east);
      \fill[UCRGold] ([yshift=-0.55cm]current page.north west) rectangle ([yshift=-0.62cm]current page.north east);
      \node[anchor=west, text=white, font=\small\bfseries] at ([xshift=5mm,yshift=-0.275cm]current page.north west)
        {ECON 282E \,$\cdot$\, Foundations of Macroeconomics};
      \node[anchor=east, text=white, font=\small] at ([xshift=-5mm,yshift=-0.275cm]current page.north east)
        {UC Riverside \,$\cdot$\, Fall 2026};
    \end{tikzpicture}
    \vspace*{1.1cm}
    {\usebeamerfont{title}\color{black!55}\large Session \decksession\ \,$\cdot$\, Deck \decksession.\deckid\par}
    \vspace{0.25cm}
    {\usebeamerfont{title}\color{RedTitle}\LARGE\bfseries \decktitletext\par}
    \vspace{0.7cm}
    {\large Zhigang Feng\par}
    \vspace{0.15cm}
    {\color{black!65}\deckdate\par}
    \vfill
    {\footnotesize\itshape\color{black!55}Build it on paper $\;\to\;$ solve it on the machine $\;\to\;$ push it to the frontier with AI\par}
  \end{frame}}

% ---------------------------------------------------------------------------
% Section divider (plain frame, not counted)
\newcommand{\sectiondivider}[2]{%
\begin{frame}[plain,noframenumbering]
\begin{center}
\vfill
{\Large\textbf{#1}\par}
\vspace{0.35cm}
{\normalsize #2\par}
\vfill
\end{center}
\end{frame}
}

% ---------------------------------------------------------------------------
% Boxes
\newtcolorbox{conceptbox}[1][]{colback=blue!3, colframe=RedTitle!55, boxrule=0.35mm,
  arc=1mm, left=2mm, right=2mm, top=1mm, bottom=1mm, #1}
\newtcolorbox{cautionbox}[1][]{colback=red!3, colframe=red!50!black, boxrule=0.35mm,
  arc=1mm, left=2mm, right=2mm, top=1mm, bottom=1mm, #1}
\newtcolorbox{examplebox}[1][]{colback=green!3, colframe=green!45!black, boxrule=0.35mm,
  arc=1mm, left=2mm, right=2mm, top=1mm, bottom=1mm, #1}
\newtcolorbox{takeawaybox}[1][]{colback=VeryLightGray, colframe=RedTitle!55, boxrule=0.35mm,
  arc=1mm, left=2mm, right=2mm, top=1mm, bottom=1mm, #1}
\newtcolorbox{researchbox}[1][]{enhanced, colback=violet!4, colframe=violet!60!black,
  boxrule=0.35mm, arc=1mm, left=2mm, right=2mm, top=1mm, bottom=1mm,
  fonttitle=\bfseries\small, title={Research in the wild}, #1}

\newcommand{\compacttables}{\renewcommand{\arraystretch}{1.15}}
\newcommand{\src}[1]{{\tiny\color{black!55}#1}}

% ---------------------------------------------------------------------------
\input{../preamble/course_map.tex}

%%   \decktitle{1}{A1}{Who I Am \& Why This Course}{September 24, 2026}
%%   \begin{document}
%%   \makedecktitle
%%   ...
%%
%% Adapted from Courses/AI-ECON-2026/source/shared_preamble.tex (Metropolis theme,
%% concept/caution/example/takeaway boxes, \sectiondivider). Changes for this course:
%%   * course title block (\decktitle / \makedecktitle) instead of \makebeamertitle
%%   * \zh{...} typesets Chinese (book titles) under pdflatex via CJKutf8
%%   * researchbox for the "Research in the wild" frames
%%   * section pages off; TikZ libraries and the course map loaded here
%% Build with pdflatex only (two passes); tools/build_slides.py does this for you.

\documentclass[10pt,english,aspectratio=169]{beamer}

\usepackage[T1]{fontenc}
\usepackage{textcomp}
\usepackage[utf8]{inputenc}
\usepackage{babel}
\usepackage{CJKutf8}

\usepackage{amsmath, amssymb, amsthm, graphicx}
\usepackage{booktabs, multirow, tabularx}
\usepackage{tikz}
\usetikzlibrary{positioning,arrows.meta,shapes,calc,fit,backgrounds}
\usepackage{pgfplots}
\pgfplotsset{compat=1.16}
\usepackage[most]{tcolorbox}
\tcbuselibrary{skins, breakable}
\usepackage{listings}
\usepackage{xcolor}

\lstset{
  basicstyle=\ttfamily\small,
  keywordstyle=\color{blue!80!black}\bfseries,
  commentstyle=\color{green!50!black}\itshape,
  stringstyle=\color{orange!80!black},
  backgroundcolor=\color{gray!8},
  frame=single,
  rulecolor=\color{gray!40},
  breaklines=true,
  showstringspaces=false,
  numbers=none,
}

\usetheme{metropolis}
\usefonttheme{professionalfonts}
\metroset{sectionpage=none}

\definecolor{LightGray}{RGB}{230,230,230}
\definecolor{RedTitle}{RGB}{0,0,200}
\definecolor{VeryLightGray}{RGB}{245,245,245}
\definecolor{DarkText}{RGB}{0,0,0}
\definecolor{UCRBlue}{RGB}{0,45,106}
\definecolor{UCRGold}{RGB}{241,171,0}

% Stage colours of the course arc (used by the course map and elsewhere)
\colorlet{StageTrunk}{blue!12}
\colorlet{StageBranch}{cyan!14}
\colorlet{StageMachine}{gray!18}
\colorlet{StageWall}{red!14}
\colorlet{StageTools}{orange!20}
\colorlet{StageData}{green!16}
\colorlet{StageAgents}{violet!14}

\hypersetup{bookmarks=false,breaklinks=true,pdfborder={0 0 0},colorlinks=true,
  linkcolor=DarkText,urlcolor=RedTitle}

\setbeamercolor{background canvas}{bg=white}
\setbeamercolor{title}{fg=RedTitle}
\setbeamercolor{frametitle}{bg=VeryLightGray,fg=DarkText}
\setbeamercolor{normal text}{fg=DarkText}
\setbeamercolor{alerted text}{fg=RedTitle}
\setbeamersize{text margin left=5mm, text margin right=5mm}
\addtobeamertemplate{frametitle}{}{\vspace{-0.5em}}

\setbeamertemplate{navigation symbols}{}
\setbeamertemplate{footline}{%
  \hspace{5mm}{\usebeamerfont{page number in head/foot}\color{black!45}%
  ECON 282E \,$\cdot$\, \insertshortdeck}\hfill%
  \usebeamercolor[fg]{page number in head/foot}%
  \usebeamerfont{page number in head/foot}%
  \insertframenumber\,/\,\inserttotalframenumber\kern1em\vskip2pt%
}

% ---------------------------------------------------------------------------
% Chinese text under pdflatex (book titles etc.)
\newcommand{\zh}[1]{\begin{CJK*}{UTF8}{gbsn}#1\end{CJK*}}

% ---------------------------------------------------------------------------
% Deck title block
%   \decktitle{session}{deck id}{title}{date}
\newcommand{\insertshortdeck}{}
\newcommand{\decksession}{}
\newcommand{\deckid}{}
\newcommand{\decktitletext}{}
\newcommand{\deckdate}{}
\newcommand{\decktitle}[4]{%
  \renewcommand{\decksession}{#1}%
  \renewcommand{\deckid}{#2}%
  \renewcommand{\decktitletext}{#3}%
  \renewcommand{\deckdate}{#4}%
  \renewcommand{\insertshortdeck}{Session #1 \,$\cdot$\, Deck #1.#2}%
  \title{#3}%
  \author{Zhigang Feng}%
  \date{#4}%
}
\newcommand{\makedecktitle}{%
  \begin{frame}[plain,noframenumbering]
    \begin{tikzpicture}[remember picture,overlay]
      \fill[UCRBlue] (current page.north west) rectangle ([yshift=-0.55cm]current page.north east);
      \fill[UCRGold] ([yshift=-0.55cm]current page.north west) rectangle ([yshift=-0.62cm]current page.north east);
      \node[anchor=west, text=white, font=\small\bfseries] at ([xshift=5mm,yshift=-0.275cm]current page.north west)
        {ECON 282E \,$\cdot$\, Foundations of Macroeconomics};
      \node[anchor=east, text=white, font=\small] at ([xshift=-5mm,yshift=-0.275cm]current page.north east)
        {UC Riverside \,$\cdot$\, Fall 2026};
    \end{tikzpicture}
    \vspace*{1.1cm}
    {\usebeamerfont{title}\color{black!55}\large Session \decksession\ \,$\cdot$\, Deck \decksession.\deckid\par}
    \vspace{0.25cm}
    {\usebeamerfont{title}\color{RedTitle}\LARGE\bfseries \decktitletext\par}
    \vspace{0.7cm}
    {\large Zhigang Feng\par}
    \vspace{0.15cm}
    {\color{black!65}\deckdate\par}
    \vfill
    {\footnotesize\itshape\color{black!55}Build it on paper $\;\to\;$ solve it on the machine $\;\to\;$ push it to the frontier with AI\par}
  \end{frame}}

% ---------------------------------------------------------------------------
% Section divider (plain frame, not counted)
\newcommand{\sectiondivider}[2]{%
\begin{frame}[plain,noframenumbering]
\begin{center}
\vfill
{\Large\textbf{#1}\par}
\vspace{0.35cm}
{\normalsize #2\par}
\vfill
\end{center}
\end{frame}
}

% ---------------------------------------------------------------------------
% Boxes
\newtcolorbox{conceptbox}[1][]{colback=blue!3, colframe=RedTitle!55, boxrule=0.35mm,
  arc=1mm, left=2mm, right=2mm, top=1mm, bottom=1mm, #1}
\newtcolorbox{cautionbox}[1][]{colback=red!3, colframe=red!50!black, boxrule=0.35mm,
  arc=1mm, left=2mm, right=2mm, top=1mm, bottom=1mm, #1}
\newtcolorbox{examplebox}[1][]{colback=green!3, colframe=green!45!black, boxrule=0.35mm,
  arc=1mm, left=2mm, right=2mm, top=1mm, bottom=1mm, #1}
\newtcolorbox{takeawaybox}[1][]{colback=VeryLightGray, colframe=RedTitle!55, boxrule=0.35mm,
  arc=1mm, left=2mm, right=2mm, top=1mm, bottom=1mm, #1}
\newtcolorbox{researchbox}[1][]{enhanced, colback=violet!4, colframe=violet!60!black,
  boxrule=0.35mm, arc=1mm, left=2mm, right=2mm, top=1mm, bottom=1mm,
  fonttitle=\bfseries\small, title={Research in the wild}, #1}

\newcommand{\compacttables}{\renewcommand{\arraystretch}{1.15}}
\newcommand{\src}[1]{{\tiny\color{black!55}#1}}

% ---------------------------------------------------------------------------
%% Course map: the recurring "you are here" slide for ECON 282E.
%% Loaded by course_preamble.tex. Pattern follows AI-ECON-2026 lec01_map.tex, but the map
%% shows the whole ten-session course rather than one lecture.
%%
%%   \CourseMap{s}{label}   s = 1..10 highlights session s and prints "you are here: label"
%%                          (e.g. \CourseMap{1}{1.A2}); earlier sessions get a check mark.
%%   \CourseMap{0}{}        full overview, nothing highlighted.
%%
%% Note: TikZ option lists are comma-split before TeX conditionals run, so every \ifnum
%% below sits outside [...] option lists.

\newcommand{\CourseMap}[2]{%
\begin{center}
\begin{tikzpicture}[
    sess/.style={rectangle, rounded corners=2pt, draw=black!45, minimum width=1.34cm,
        minimum height=1.12cm, text width=1.26cm, align=center, inner sep=1pt},
    stage/.style={font=\tiny\bfseries, text=black!70},
]
  \foreach \i/\d/\t/\c in {%
      1/Sep 24/Intro \\ Growth/StageTrunk,
      2/Oct 1/HA $\cdot$ OLG \\ Policy/StageBranch,
      3/Oct 8/Num.\ DP \\ Python/StageMachine,
      4/Oct 15/Numerics \\ I/StageMachine,
      5/Oct 22/Numerics \\ II/StageMachine,
      6/Oct 29/The wall \\ \textbf{Midterm}/StageWall,
      7/Nov 5/ML \\ Deep nets/StageTools,
      8/Nov 12/RL \\ HA $+$ DL/StageTools,
      9/Nov 19/Text \\ LLMs/StageData,
      10/Dec 3/Agents \\ \textbf{Projects}/StageAgents} {
    \node[sess, fill=\c] (n\i) at ({1.46*(\i-1)}, 0)
      {{\scriptsize\bfseries S\i}\\[-1pt]{\tiny\color{black!60}\d}\\[1pt]{\tiny \t}};
    \ifnum\i<#1
      \node[font=\scriptsize, text=green!45!black] at ([xshift=-2.5mm,yshift=-2.2mm]n\i.north east) {$\checkmark$};
    \fi
    \ifnum\i=#1
      \draw[RedTitle, very thick, rounded corners=2pt]
        ([xshift=-0.6mm,yshift=-0.6mm]n\i.south west) rectangle ([xshift=0.6mm,yshift=0.6mm]n\i.north east);
      % keep the label inside the picture at both ends of the row
      \ifnum\i<3
        \node[font=\scriptsize\bfseries, text=RedTitle, anchor=south west] at ([yshift=1.2mm]n\i.north west)
          {$\blacktriangledown$\,you are here: #2};
      \else\ifnum\i>8
        \node[font=\scriptsize\bfseries, text=RedTitle, anchor=south east] at ([yshift=1.2mm]n\i.north east)
          {you are here: #2\,$\blacktriangledown$};
      \else
        \node[font=\scriptsize\bfseries, text=RedTitle, anchor=south] at ([yshift=1.2mm]n\i.north)
          {you are here: #2\,$\blacktriangledown$};
      \fi\fi
    \fi
  }
  % stage band under the sessions
  \foreach \a/\b/\lab/\c in {1/1/trunk/StageTrunk, 2/2/branches/StageBranch,
      3/5/the machine $\cdot$ classical methods/StageMachine, 6/6/the wall/StageWall,
      7/8/new tools/StageTools, 9/9/new data/StageData, 10/10/colleagues/StageAgents} {
    \fill[\c] ([yshift=-1.5mm]n\a.south west) rectangle ([yshift=-4.6mm]n\b.south east);
    \path ([yshift=-1.5mm]n\a.south west) -- ([yshift=-4.6mm]n\b.south east)
      node[midway, stage] {\lab};
  }
  \node[font=\footnotesize\itshape, text=black!70] at ([yshift=-9.5mm]$(n1.south)!0.5!(n10.south)$)
    {Build it on paper $\;\to\;$ solve it on the machine $\;\to\;$ push it to the frontier with AI};
\end{tikzpicture}
\end{center}
}


%%   \decktitle{1}{A1}{Who I Am \& Why This Course}{September 24, 2026}
%%   \begin{document}
%%   \makedecktitle
%%   ...
%%
%% Adapted from Courses/AI-ECON-2026/source/shared_preamble.tex (Metropolis theme,
%% concept/caution/example/takeaway boxes, \sectiondivider). Changes for this course:
%%   * course title block (\decktitle / \makedecktitle) instead of \makebeamertitle
%%   * \zh{...} typesets Chinese (book titles) under pdflatex via CJKutf8
%%   * researchbox for the "Research in the wild" frames
%%   * section pages off; TikZ libraries and the course map loaded here
%% Build with pdflatex only (two passes); tools/build_slides.py does this for you.

\documentclass[10pt,english,aspectratio=169]{beamer}

\usepackage[T1]{fontenc}
\usepackage{textcomp}
\usepackage[utf8]{inputenc}
\usepackage{babel}
\usepackage{CJKutf8}

\usepackage{amsmath, amssymb, amsthm, graphicx}
\usepackage{booktabs, multirow, tabularx}
\usepackage{tikz}
\usetikzlibrary{positioning,arrows.meta,shapes,calc,fit,backgrounds}
\usepackage{pgfplots}
\pgfplotsset{compat=1.16}
\usepackage[most]{tcolorbox}
\tcbuselibrary{skins, breakable}
\usepackage{listings}
\usepackage{xcolor}

\lstset{
  basicstyle=\ttfamily\small,
  keywordstyle=\color{blue!80!black}\bfseries,
  commentstyle=\color{green!50!black}\itshape,
  stringstyle=\color{orange!80!black},
  backgroundcolor=\color{gray!8},
  frame=single,
  rulecolor=\color{gray!40},
  breaklines=true,
  showstringspaces=false,
  numbers=none,
}

\usetheme{metropolis}
\usefonttheme{professionalfonts}
\metroset{sectionpage=none}

\definecolor{LightGray}{RGB}{230,230,230}
\definecolor{RedTitle}{RGB}{0,0,200}
\definecolor{VeryLightGray}{RGB}{245,245,245}
\definecolor{DarkText}{RGB}{0,0,0}
\definecolor{UCRBlue}{RGB}{0,45,106}
\definecolor{UCRGold}{RGB}{241,171,0}

% Stage colours of the course arc (used by the course map and elsewhere)
\colorlet{StageTrunk}{blue!12}
\colorlet{StageBranch}{cyan!14}
\colorlet{StageMachine}{gray!18}
\colorlet{StageWall}{red!14}
\colorlet{StageTools}{orange!20}
\colorlet{StageData}{green!16}
\colorlet{StageAgents}{violet!14}

\hypersetup{bookmarks=false,breaklinks=true,pdfborder={0 0 0},colorlinks=true,
  linkcolor=DarkText,urlcolor=RedTitle}

\setbeamercolor{background canvas}{bg=white}
\setbeamercolor{title}{fg=RedTitle}
\setbeamercolor{frametitle}{bg=VeryLightGray,fg=DarkText}
\setbeamercolor{normal text}{fg=DarkText}
\setbeamercolor{alerted text}{fg=RedTitle}
\setbeamersize{text margin left=5mm, text margin right=5mm}
\addtobeamertemplate{frametitle}{}{\vspace{-0.5em}}

\setbeamertemplate{navigation symbols}{}
\setbeamertemplate{footline}{%
  \hspace{5mm}{\usebeamerfont{page number in head/foot}\color{black!45}%
  ECON 282E \,$\cdot$\, \insertshortdeck}\hfill%
  \usebeamercolor[fg]{page number in head/foot}%
  \usebeamerfont{page number in head/foot}%
  \insertframenumber\,/\,\inserttotalframenumber\kern1em\vskip2pt%
}

% ---------------------------------------------------------------------------
% Chinese text under pdflatex (book titles etc.)
\newcommand{\zh}[1]{\begin{CJK*}{UTF8}{gbsn}#1\end{CJK*}}

% ---------------------------------------------------------------------------
% Deck title block
%   \decktitle{session}{deck id}{title}{date}
\newcommand{\insertshortdeck}{}
\newcommand{\decksession}{}
\newcommand{\deckid}{}
\newcommand{\decktitletext}{}
\newcommand{\deckdate}{}
\newcommand{\decktitle}[4]{%
  \renewcommand{\decksession}{#1}%
  \renewcommand{\deckid}{#2}%
  \renewcommand{\decktitletext}{#3}%
  \renewcommand{\deckdate}{#4}%
  \renewcommand{\insertshortdeck}{Session #1 \,$\cdot$\, Deck #1.#2}%
  \title{#3}%
  \author{Zhigang Feng}%
  \date{#4}%
}
\newcommand{\makedecktitle}{%
  \begin{frame}[plain,noframenumbering]
    \begin{tikzpicture}[remember picture,overlay]
      \fill[UCRBlue] (current page.north west) rectangle ([yshift=-0.55cm]current page.north east);
      \fill[UCRGold] ([yshift=-0.55cm]current page.north west) rectangle ([yshift=-0.62cm]current page.north east);
      \node[anchor=west, text=white, font=\small\bfseries] at ([xshift=5mm,yshift=-0.275cm]current page.north west)
        {ECON 282E \,$\cdot$\, Foundations of Macroeconomics};
      \node[anchor=east, text=white, font=\small] at ([xshift=-5mm,yshift=-0.275cm]current page.north east)
        {UC Riverside \,$\cdot$\, Fall 2026};
    \end{tikzpicture}
    \vspace*{1.1cm}
    {\usebeamerfont{title}\color{black!55}\large Session \decksession\ \,$\cdot$\, Deck \decksession.\deckid\par}
    \vspace{0.25cm}
    {\usebeamerfont{title}\color{RedTitle}\LARGE\bfseries \decktitletext\par}
    \vspace{0.7cm}
    {\large Zhigang Feng\par}
    \vspace{0.15cm}
    {\color{black!65}\deckdate\par}
    \vfill
    {\footnotesize\itshape\color{black!55}Build it on paper $\;\to\;$ solve it on the machine $\;\to\;$ push it to the frontier with AI\par}
  \end{frame}}

% ---------------------------------------------------------------------------
% Section divider (plain frame, not counted)
\newcommand{\sectiondivider}[2]{%
\begin{frame}[plain,noframenumbering]
\begin{center}
\vfill
{\Large\textbf{#1}\par}
\vspace{0.35cm}
{\normalsize #2\par}
\vfill
\end{center}
\end{frame}
}

% ---------------------------------------------------------------------------
% Boxes
\newtcolorbox{conceptbox}[1][]{colback=blue!3, colframe=RedTitle!55, boxrule=0.35mm,
  arc=1mm, left=2mm, right=2mm, top=1mm, bottom=1mm, #1}
\newtcolorbox{cautionbox}[1][]{colback=red!3, colframe=red!50!black, boxrule=0.35mm,
  arc=1mm, left=2mm, right=2mm, top=1mm, bottom=1mm, #1}
\newtcolorbox{examplebox}[1][]{colback=green!3, colframe=green!45!black, boxrule=0.35mm,
  arc=1mm, left=2mm, right=2mm, top=1mm, bottom=1mm, #1}
\newtcolorbox{takeawaybox}[1][]{colback=VeryLightGray, colframe=RedTitle!55, boxrule=0.35mm,
  arc=1mm, left=2mm, right=2mm, top=1mm, bottom=1mm, #1}
\newtcolorbox{researchbox}[1][]{enhanced, colback=violet!4, colframe=violet!60!black,
  boxrule=0.35mm, arc=1mm, left=2mm, right=2mm, top=1mm, bottom=1mm,
  fonttitle=\bfseries\small, title={Research in the wild}, #1}

\newcommand{\compacttables}{\renewcommand{\arraystretch}{1.15}}
\newcommand{\src}[1]{{\tiny\color{black!55}#1}}

% ---------------------------------------------------------------------------
%% Course map: the recurring "you are here" slide for ECON 282E.
%% Loaded by course_preamble.tex. Pattern follows AI-ECON-2026 lec01_map.tex, but the map
%% shows the whole ten-session course rather than one lecture.
%%
%%   \CourseMap{s}{label}   s = 1..10 highlights session s and prints "you are here: label"
%%                          (e.g. \CourseMap{1}{1.A2}); earlier sessions get a check mark.
%%   \CourseMap{0}{}        full overview, nothing highlighted.
%%
%% Note: TikZ option lists are comma-split before TeX conditionals run, so every \ifnum
%% below sits outside [...] option lists.

\newcommand{\CourseMap}[2]{%
\begin{center}
\begin{tikzpicture}[
    sess/.style={rectangle, rounded corners=2pt, draw=black!45, minimum width=1.34cm,
        minimum height=1.12cm, text width=1.26cm, align=center, inner sep=1pt},
    stage/.style={font=\tiny\bfseries, text=black!70},
]
  \foreach \i/\d/\t/\c in {%
      1/Sep 24/Intro \\ Growth/StageTrunk,
      2/Oct 1/HA $\cdot$ OLG \\ Policy/StageBranch,
      3/Oct 8/Num.\ DP \\ Python/StageMachine,
      4/Oct 15/Numerics \\ I/StageMachine,
      5/Oct 22/Numerics \\ II/StageMachine,
      6/Oct 29/The wall \\ \textbf{Midterm}/StageWall,
      7/Nov 5/ML \\ Deep nets/StageTools,
      8/Nov 12/RL \\ HA $+$ DL/StageTools,
      9/Nov 19/Text \\ LLMs/StageData,
      10/Dec 3/Agents \\ \textbf{Projects}/StageAgents} {
    \node[sess, fill=\c] (n\i) at ({1.46*(\i-1)}, 0)
      {{\scriptsize\bfseries S\i}\\[-1pt]{\tiny\color{black!60}\d}\\[1pt]{\tiny \t}};
    \ifnum\i<#1
      \node[font=\scriptsize, text=green!45!black] at ([xshift=-2.5mm,yshift=-2.2mm]n\i.north east) {$\checkmark$};
    \fi
    \ifnum\i=#1
      \draw[RedTitle, very thick, rounded corners=2pt]
        ([xshift=-0.6mm,yshift=-0.6mm]n\i.south west) rectangle ([xshift=0.6mm,yshift=0.6mm]n\i.north east);
      % keep the label inside the picture at both ends of the row
      \ifnum\i<3
        \node[font=\scriptsize\bfseries, text=RedTitle, anchor=south west] at ([yshift=1.2mm]n\i.north west)
          {$\blacktriangledown$\,you are here: #2};
      \else\ifnum\i>8
        \node[font=\scriptsize\bfseries, text=RedTitle, anchor=south east] at ([yshift=1.2mm]n\i.north east)
          {you are here: #2\,$\blacktriangledown$};
      \else
        \node[font=\scriptsize\bfseries, text=RedTitle, anchor=south] at ([yshift=1.2mm]n\i.north)
          {you are here: #2\,$\blacktriangledown$};
      \fi\fi
    \fi
  }
  % stage band under the sessions
  \foreach \a/\b/\lab/\c in {1/1/trunk/StageTrunk, 2/2/branches/StageBranch,
      3/5/the machine $\cdot$ classical methods/StageMachine, 6/6/the wall/StageWall,
      7/8/new tools/StageTools, 9/9/new data/StageData, 10/10/colleagues/StageAgents} {
    \fill[\c] ([yshift=-1.5mm]n\a.south west) rectangle ([yshift=-4.6mm]n\b.south east);
    \path ([yshift=-1.5mm]n\a.south west) -- ([yshift=-4.6mm]n\b.south east)
      node[midway, stage] {\lab};
  }
  \node[font=\footnotesize\itshape, text=black!70] at ([yshift=-9.5mm]$(n1.south)!0.5!(n10.south)$)
    {Build it on paper $\;\to\;$ solve it on the machine $\;\to\;$ push it to the frontier with AI};
\end{tikzpicture}
\end{center}
}


\graphicspath{{./figures/}}

\decktitle{7}{B2}{Deep VFI and Actor--Critic}{Thursday, November 5, 2026}

\begin{document}

\makedecktitle

%=============================================================================
\begin{frame}{Where We Are}
\CourseMap{7}{7.B2}
\vspace{-1mm}
\begin{center}
\small \textbf{Block B:} 7.B1 the model as a loss function $\;\cdot\;$ \textbf{7.B2} deep VFI and actor--critic $\;\cdot\;$ 7.B3 structure and validation.
\end{center}
\end{frame}

%=============================================================================
\begin{frame}{The Other Road}
\begin{columns}[T]
\begin{column}{0.54\textwidth}
\small
7.B1 required first-order conditions --- a real restriction, and it bites exactly where macroeconomics gets interesting:
\smallskip
\begin{itemize}\small\setlength{\itemsep}{0.7mm}
  \item \textbf{occasionally binding constraints} --- the Euler equation holds with an inequality;
  \item \textbf{discrete choice} --- default, retirement, entry: no derivative to set to zero;
  \item \textbf{kinks} --- a borrowing limit, an irreversibility.
\end{itemize}
\smallskip

So go back one step, to the object the FOCs were derived \emph{from}:
\[
  V(k)=\max_{k'}\bigl\{u(c)+\beta V(k')\bigr\},
\]
also a functional equation with a residual. \textbf{Make the residual small and you have solved it} --- and \textbf{no derivatives of the model are needed}, only derivatives of the \emph{network}. The price is the $\max$, which now sits \emph{inside} the loss.
\end{column}
\begin{column}{0.43\textwidth}
\begin{researchbox}\small
\textbf{The route with theory attached.}
\medskip
Feng \& Santos construct Markov equilibrium for a time-inconsistent economy from \textbf{value-function operators}, and obtain existence \emph{constructively} --- as the limit of a sequence of $\varepsilon$-equilibria.
\medskip
Where an economy is hard enough that \textbf{existence is itself in question}, the value function is where the argument gets made. That is a reason to prefer this road beyond convenience.
\end{researchbox}
\end{column}
\end{columns}
\vspace{1mm}
\src{Feng \& Santos, ``Markov Equilibria with Quasi-Hyperbolic Discounting'', working paper.}
\end{frame}

%=============================================================================
\begin{frame}{Replace the Grid with a Network}
\begin{columns}[T]
\begin{column}{0.53\textwidth}
\small
Instead of storing $V(k)$ at grid points, approximate both unknowns:
\[
  V(k)\approx\mathcal{V}(k;\theta_v),\qquad g(k)\approx\mathcal{G}(k;\theta_g).
\]
\medskip

The loop keeps classical VFI's rhythm exactly --- \emph{evaluate the current policy, then improve it} --- with two substitutions:
\begin{enumerate}\small\setlength{\itemsep}{0.9mm}
  \item \textbf{sample} states $k\sim D(k)$ instead of laying down a grid;
  \item \textbf{fit} the networks by gradient descent instead of writing into a table.
\end{enumerate}
\end{column}
\begin{column}{0.44\textwidth}
\begin{takeawaybox}[title=\textbf{Why it scales}]\small
A network handles high-dimensional $k$ without a mesh: states are sampled, the network \textbf{generalises} between them, and the parameter count grows polynomially in $d$ rather than as $N^{d}$.
\end{takeawaybox}
\medskip
\begin{examplebox}\footnotesize
``Generalises between them'' is doing the work, and it is also the risk. A grid makes no claim about the points it does not contain. \textbf{A network always answers}, everywhere, whether or not it should.
\end{examplebox}
\end{column}
\end{columns}
\end{frame}

%=============================================================================
\begin{frame}{Deep VFI, the Whole Loop in One Place}
\begin{conceptbox}[title=\textbf{Deep value function iteration}]\footnotesize
Represent $V\approx\mathcal{V}(\cdot;\theta_v)$ and $g\approx\mathcal{G}(\cdot;\theta_g)$, and iterate:
\begin{enumerate}\footnotesize\setlength{\itemsep}{0.8mm}
  \item \textbf{Sample} a batch of states $\{k_i\}\sim D(k)$ --- no grid.
  \item \textbf{Value step (the critic).} Fit $\mathcal{V}$ by least squares to the Bellman target
  $\hat V(k)=u\bigl(k,g(k)\bigr)+\beta\,\mathcal{V}\bigl(f(k,g(k))\bigr)$.
  \item \textbf{Policy step (the actor).} Improve $\mathcal{G}$ by gradient \emph{ascent} on
  $u\bigl(k,g(k)\bigr)+\beta\,\mathcal{V}\bigl(f(k,g(k))\bigr)$.
  \item \textbf{Repeat} until both networks stop moving.
\end{enumerate}
\end{conceptbox}
\medskip
\begin{columns}[T]
\begin{column}{0.49\textwidth}
\begin{takeawaybox}\footnotesize
Steps 2 and 3 are \textbf{policy evaluation} and \textbf{policy improvement} --- 1.B2's two halves of value-function iteration, unchanged in logic and wholly changed in machinery.
\end{takeawaybox}
\end{column}
\begin{column}{0.49\textwidth}
\begin{cautionbox}\footnotesize
Notice what is \emph{not} in step 3: an $\arg\max$. The policy network is nudged in the direction that raises the objective, rather than being resolved at each state.
\end{cautionbox}
\end{column}
\end{columns}
\end{frame}

%=============================================================================
\begin{frame}{Deep VFI Against Classical VFI}
\vspace{-1mm}
\begin{center}\footnotesize
\begin{tabularx}{\linewidth}{@{}>{\raggedright\arraybackslash}p{2.6cm} >{\raggedright\arraybackslash}X >{\raggedright\arraybackslash}X@{}}
\toprule
& \textbf{Classical VFI (3.A2)} & \textbf{Deep VFI} \\
\midrule
state space & a fixed grid $\{k_1,\dots,k_N\}$ & a sampled batch $k_i\sim D(k)$ \\
representation & table, plus interpolation & a network $\mathcal{V}(k;\theta_v)$ \\
value update & exact $\max_{k'}$ at each point & least-squares fit to $\hat V$ \\
policy & $\arg\max$ by grid search & a learned network, by ascent \\
expectation & quadrature on a grid & exact sum, or Monte Carlo \\
cost in $d$ & $N^{d}$ & polynomial \\
\textbf{convergence} & \textbf{a contraction (Banach)} & \textbf{no general guarantee} \\
\bottomrule
\end{tabularx}
\end{center}
\medskip
\begin{columns}[T]
\begin{column}{0.49\textwidth}
\begin{takeawaybox}\footnotesize
\textbf{The trade, in one line.} Give up the exact maximisation and the contraction guarantee; gain dimension, and the ability to fit scattered, simulated data.
\end{takeawaybox}
\end{column}
\begin{column}{0.49\textwidth}
\begin{cautionbox}\footnotesize
\textbf{Read the last row slowly.} 3.A2's error bound $\lVert V_n-V^{*}\rVert\le\beta^{n}\lVert V_0-V^{*}\rVert/(1-\beta)$ told you \emph{how far you were} without knowing the answer. \textbf{Nothing here replaces it.}
\end{cautionbox}
\end{column}
\end{columns}
\vspace{1mm}
\src{Losing the bound is why 7.B3 exists.}
\end{frame}

%=============================================================================
\begin{frame}{The Value Step, and the Detach That Makes It Work}
\begin{columns}[T]
\begin{column}{0.55\textwidth}
\small
Hold the policy $g^{(n-1)}$ fixed. The one-step Bellman target is
\[
  \hat V(k)=u\bigl(k,g^{(n-1)}(k)\bigr)+\beta\,\mathcal{V}^{(n-1)}\bigl(f(k,g^{(n-1)}(k))\bigr),
\]
and the value network is fitted to it by least squares:
\[
  \min_{\theta_v}\;\mathbb{E}_{k\sim D}\Bigl[\bigl(\mathcal{V}(k;\theta_v)-\hat V(k)\bigr)^{2}\Bigr].
\]
\medskip

\textbf{This is supervised learning again}, with the ``label'' at each sampled state being the Bellman target.
\end{column}
\begin{column}{0.42\textwidth}
\begin{cautionbox}[title=\textbf{\texttt{detach()} the target}]\small
$\hat V$ contains $\mathcal{V}^{(n-1)}$, which is the \emph{same network}. If the gradient is allowed to flow into the target as well as the prediction, the optimizer can lower the loss by \textbf{moving the target} rather than fitting it.
\medskip
Detaching freezes the target for the step. It is one line of code and the difference between converging and not.
\end{cautionbox}
\medskip
\begin{examplebox}\footnotesize
The idea has a name in deep RL --- the \textbf{target network} --- and it returns in S8.
\end{examplebox}
\end{column}
\end{columns}
\end{frame}

%=============================================================================
\begin{frame}{The Policy Step, and Why Two Networks}
\begin{columns}[T]
\begin{column}{0.51\textwidth}
\small
\textbf{Where is the maximisation actually solved?}
\smallskip

\begin{center}\footnotesize
\begin{tabularx}{\linewidth}{@{}>{\raggedright\arraybackslash}p{2.2cm} >{\raggedright\arraybackslash}X@{}}
\toprule
grid VFI (3.A2) & at every grid point \\
projection (5.B6) & at the \textbf{collocation nodes} only \\
deep VFI & at every \textbf{sampled} state, every batch \\
\bottomrule
\end{tabularx}
\end{center}
\smallskip

\textbf{What makes it hard is not that a method is classical}, but whether the approximant \textbf{preserves concavity}. A concave interpolant leaves a concave program with one maximiser, which 4.A2's golden section finds every time; a rippling polynomial --- or a generic network --- does not.
\smallskip

So let a second network \textbf{output the action directly}, $k'=\mathcal{G}(k;\theta_g)$, improved by gradient \emph{ascent} on $u(k,\mathcal{G}(k))+\beta\,\mathcal{V}(f(k,\mathcal{G}(k)))$.
\end{column}
\begin{column}{0.46\textwidth}
\begin{conceptbox}[title=\textbf{The two roles}]\footnotesize
\textbf{The actor} $\mathcal{G}$ --- the doer. State in, action out; it never maximises anything, it is \emph{improved}.
\medskip
\textbf{The critic} $\mathcal{V}$ --- the judge. State in, value out; it scores what the actor did. The actor proposes, the critic evaluates, both improve together.
\end{conceptbox}
\medskip
\begin{takeawaybox}\footnotesize
\textbf{Actor--critic does not avoid the maximisation --- it \emph{amortises} it.} The actor learns the $\arg\max$ as a \emph{function of the state}: one gradient step over parameters instead of one exact maximisation everywhere.
\medskip
Which is why 7.B3 matters: \textbf{restore concavity and the inner problem is well behaved again.}
\end{takeawaybox}
\end{column}
\end{columns}
\end{frame}

%=============================================================================
\begin{frame}{Two Moving Targets}
\begin{columns}[T]
\begin{column}{0.53\textwidth}
\small
The awkwardness of the scheme is that \textbf{neither network is chasing a fixed object}.
\medskip
\begin{itemize}\small\setlength{\itemsep}{1.1mm}
  \item The critic is fitted to a target built from the \emph{current} policy. Improve the policy and the target moves.
  \item The actor ascends an objective built from the \emph{current} critic. Refit the critic and the landscape moves.
\end{itemize}
\medskip

Classical VFI has this too --- but there the value step is a \textbf{contraction}, which is precisely what makes the alternation converge. Replace the exact update with a least-squares fit on sampled states and the contraction is gone.
\end{column}
\begin{column}{0.44\textwidth}
\begin{takeawaybox}[title=\textbf{What is done in practice}]\small
\textbf{Freeze and sync.} Hold a frozen copy of the target network and refresh it every $n$ iterations, so each side chases something stationary for a while.
\medskip
\textbf{Unequal step sizes.} Let the critic learn faster than the actor, so the judge is roughly right before the doer trusts it.
\end{takeawaybox}
\medskip
\begin{examplebox}\footnotesize
Both devices reappear in S8's frontier solver, where the frozen copy is what imposes \emph{perceived $=$ actual} in equilibrium.
\end{examplebox}
\end{column}
\end{columns}
\end{frame}

%=============================================================================
\begin{frame}{One Step or Many: a Bias--Variance Choice}
\begin{columns}[T]
\begin{column}{0.53\textwidth}
\small
\textbf{One step.} $\hat V(k_i)=u\bigl(k_i,g(k_i)\bigr)+\beta\,\mathcal{V}(k_{i,1})$.
\smallskip

\emph{Low variance} --- one transition. \emph{High bias} --- it leans almost entirely on the current $\mathcal{V}$, which is wrong early on.
\smallskip

\textbf{$T$ steps.}
\[
  \hat V(k_i)=\sum_{t=0}^{T-1}\beta^{t}u\bigl(k_{i,t},g(k_{i,t})\bigr)+\beta^{T}\mathcal{V}(k_{i,T}).
\]
\emph{Lower bias} --- $T$ periods of real simulated reward, the network used only for the tail. \emph{Higher variance} --- paths differ.
\end{column}
\begin{column}{0.44\textwidth}
\begin{takeawaybox}[title=\textbf{The dial}]\small
$T$ trades the network's error against sampling noise. Typical values are $30$--$50$: long enough for $\beta^{T}$ to make the tail nearly irrelevant, short enough that the paths have not dispersed.
\end{takeawaybox}
\smallskip
{\footnotesize At $\beta=0.99$ quarterly, $\beta^{50}\approx0.61$ and $\beta^{200}\approx0.13$: \textbf{the discount factor sets how much horizon must be simulated} before the critic's error stops mattering.
\medskip
S8 names the family --- $T=1$ is temporal difference, $T=\infty$ is Monte Carlo, in between is TD($\lambda$).}
\end{column}
\end{columns}
\end{frame}

%=============================================================================
\begin{frame}{Where to Sample, and Why Not Chebyshev}
\begin{columns}[T]
\begin{column}{0.53\textwidth}
\small
Chebyshev and finite elements must approximate over the \textbf{entire} domain --- that is what a basis expansion is. As $d$ grows, covering the whole domain is exactly what becomes impossible.
\medskip

Sampling from the model's own \textbf{stationary distribution} does something different in kind: it spends the approximation budget on the region the economy actually visits, and ignores the rest.
\medskip

That is the synergy. \textbf{Monte Carlo generates data where it matters; the network fits scattered, mesh-free data.} Neither half works without the other.
\end{column}
\begin{column}{0.44\textwidth}
\begin{cautionbox}[title=\textbf{Three concerns, and their remedies}]\small
\textbf{Rare but important states.} A crisis is off the ergodic path by construction --- mix in \emph{targeted} or importance sampling.
\medskip
\textbf{Multiple steady states.} One stationary distribution can miss an entire basin.
\medskip
\textbf{Out-of-distribution behaviour.} The network answers confidently outside the region it saw. \textbf{Stress-test separately.}
\end{cautionbox}
\medskip
\begin{examplebox}\footnotesize
7.B1 measured the first of these: two decades of accuracy lost at the edge of the sampled band.
\end{examplebox}
\end{column}
\end{columns}
\end{frame}

%=============================================================================
\begin{frame}{Equilibrium Conditions Can Be Built In}
\begin{columns}[T]
\begin{column}{0.51\textwidth}
\small
A network's output can be made to satisfy economic restrictions rather than merely approach them. Four devices, in increasing order of strength:
\medskip
\begin{enumerate}\small\setlength{\itemsep}{1mm}
  \item \textbf{Output transformation} --- a sigmoid bounds the action into the feasible set. 7.B1 used it.
  \item \textbf{Penalty} --- add $\lambda\sum(\text{violation})^{2}$ to the loss.
  \item \textbf{Lagrangian} --- carry dual variables and update them during training.
  \item \textbf{Architecture} --- structure the network so an identity holds automatically.
\end{enumerate}
\end{column}
\begin{column}{0.46\textwidth}
\begin{takeawaybox}[title=\textbf{Which to reach for}]\small
\textbf{An identity} --- market clearing, a budget constraint --- should be imposed by (1) or (4). A quantity that must equal another should never be \emph{learned}.
\medskip
\textbf{An inequality} that binds sometimes needs (2) or (3), or a Fischer--Burmeister reformulation.
\end{takeawaybox}
\medskip
\begin{cautionbox}\footnotesize
A penalty coefficient is a real hyperparameter with real consequences: too small and the constraint is violated, too large and it dominates the residual and nothing else is learned.
\end{cautionbox}
\end{column}
\end{columns}
\vspace{1mm}
\src{Device (4) is 7.B3.}
\end{frame}

%=============================================================================
\begin{frame}{Macroeconomics in Reinforcement-Learning Vocabulary}
\vspace{-1mm}
\begin{center}\footnotesize
\begin{tabularx}{\linewidth}{@{}>{\raggedright\arraybackslash}p{2.4cm} >{\raggedright\arraybackslash}p{3.4cm} >{\raggedright\arraybackslash}X@{}}
\toprule
\textbf{RL} & \textbf{Economics} & \textbf{Note} \\
\midrule
agent & household or firm & the optimiser \\
environment & prices and constraints & the rules of the game \\
state $S$ & $(k,z,\Gamma)$ & wealth, shocks, \textbf{the distribution} \\
action $A$ & $c$, $i$, $l$ & the choice \\
transition $P$ & $k'=(1-\delta)k+i$ & \textbf{micro: known exactly} \\
 & $\Gamma'=H(\Gamma)$ & \textbf{macro: often unknown} \\
reward $R$ & $u(c)$ & period utility \\
discount $\gamma$ & $\beta$ & time preference \\
\bottomrule
\end{tabularx}
\end{center}
\medskip
\begin{columns}[T]
\begin{column}{0.49\textwidth}
\begin{takeawaybox}\footnotesize
\textbf{The two transition rows are the whole story of S8.} We know the household's law of motion exactly, which is why nothing in this deck needed exploration.
\end{takeawaybox}
\end{column}
\begin{column}{0.49\textwidth}
\begin{cautionbox}\footnotesize
We do \emph{not} know $H$. In Krusell--Smith it is guessed and regressed; the frontier replaces the guess with a network that takes $\Gamma$ itself as an input.
\end{cautionbox}
\end{column}
\end{columns}
\vspace{1mm}
\src{2.B1 promised this frame.}
\end{frame}

%=============================================================================
\begin{frame}{Two Practical Problems, and Their Standard Answers}
\begin{columns}[T]
\begin{column}{0.49\textwidth}
\small
\textbf{The horizon is infinite; a simulation is not.}
\medskip

Truncate at $T$ and close with the critic:
\[
  \sum_{t=0}^{T-1}\beta^{t}u_t+\beta^{T}\mathcal{V}(k_T).
\]
The tail is not dropped --- it is \emph{estimated}, and $\beta^{T}$ bounds the damage. The critic exists for exactly this.
\end{column}
\begin{column}{0.49\textwidth}
\small
\textbf{The expectation is unknown in closed form.}
\medskip

Simulate. Draw shock paths, average the realised returns, and let the noise be absorbed by the same stochastic descent that already tolerates mini-batch noise.
\medskip

Where the shock is a finite Markov chain --- as in 7.B1 --- \textbf{take the exact sum instead}. It costs one matrix product and removes a whole source of error.
\end{column}
\end{columns}
\medskip
\begin{takeawaybox}\footnotesize
\textbf{A rule worth carrying:} never sample what you can sum. The inner expectation is contamination of the residual itself, and unlike outer sampling noise it buys you nothing.
\end{takeawaybox}
\end{frame}

%=============================================================================
\begin{frame}{Euler or Bellman? A Decision, Not a Doctrine}
\vspace{-1mm}
\begin{center}\footnotesize
\begin{tabularx}{\linewidth}{@{}>{\raggedright\arraybackslash}p{2.7cm} >{\raggedright\arraybackslash}X >{\raggedright\arraybackslash}X@{}}
\toprule
& \textbf{Euler residual (7.B1)} & \textbf{Deep VFI / actor--critic (7.B2)} \\
\midrule
networks & one & two \\
requires & first-order conditions & only the Bellman equation \\
flavour & supervised & reinforcement \\
gives you & the policy & the policy \textbf{and} the value \\
handles kinks & badly & naturally \\
discrete choice & no & yes \\
accuracy when FOCs exist & \textbf{higher} & lower \\
\bottomrule
\end{tabularx}
\end{center}
\medskip
\begin{columns}[T]
\begin{column}{0.49\textwidth}
\begin{takeawaybox}\footnotesize
\textbf{Use the Euler road when the FOCs are clean}: one network, a sharper residual, and a quantity you can audit directly against a classical solve.
\end{takeawaybox}
\end{column}
\begin{column}{0.49\textwidth}
\begin{cautionbox}\footnotesize
\textbf{Use the Bellman road when they are not} --- and accept that you have given up the cleanest validation instrument along with them. 7.B3 supplies a replacement.
\end{cautionbox}
\end{column}
\end{columns}
\vspace{1mm}
\src{The value function is worth having in its own right --- welfare comparisons need it, and Euler methods do not produce one.}
\end{frame}

%=============================================================================
\begin{frame}{Two Questions This Deck Cannot Answer}
\begin{columns}[T]
\begin{column}{0.49\textwidth}
\begin{conceptbox}[title=\textbf{Why sample along the state space at all?}]\footnotesize
Because accuracy is only worth buying where the economy actually goes. But the ergodic distribution \textbf{depends on the policy being learned} --- so the sampler and the solution are a \emph{joint fixed point}, and each is chasing the other.
\end{conceptbox}
\smallskip
{\footnotesize Train only where the current policy goes and it can settle in a corner, satisfy the Bellman equation beautifully there, and be wrong everywhere else --- with a small loss reporting success. \textbf{That is on-policy against off-policy.}}
\end{column}
\begin{column}{0.49\textwidth}
\begin{conceptbox}[title=\textbf{What goes wrong in a $T$-step rollout?}]\small
\textbf{Uncertainty unfolds.} Every simulated step multiplies the branches the path could have taken, so the variance of the return estimate \textbf{compounds with $T$} --- while the bias from the critic falls.
\end{conceptbox}
\smallskip
{\footnotesize Longer horizons buy less bias and pay in variance. \textbf{Controlling that exchange --- baselines, temporal-difference targets, variance reduction --- is most of what reinforcement learning is.}}
\end{column}
\end{columns}
\smallskip
\begin{takeawaybox}\footnotesize
Both have the same shape: \textbf{the object being estimated depends on the policy doing the estimating.} Nothing here resolves that; the methods above work anyway because the growth model is forgiving. \textbf{Session 8 is where it stops being forgiving.}
\end{takeawaybox}
\end{frame}

%=============================================================================
\begin{frame}{Takeaways, and What 7.B3 Has to Fix}
\begin{columns}[T]
\begin{column}{0.55\textwidth}
\begin{takeawaybox}\small
\begin{itemize}\small\setlength{\itemsep}{1.1mm}
  \item The Bellman equation is a functional equation with a residual, so it is a loss --- and it needs \textbf{no first-order conditions}.
  \item The inner $\max$ is the obstacle, and \textbf{a second network removes it}: the actor acts, the critic judges.
  \item The value target must be \textbf{detached}, or the optimizer lowers the loss by moving the target.
  \item Sampling from the ergodic region is what makes this scale --- and is exactly why the answer \textbf{outside} that region is unconstrained.
  \item \textbf{The contraction is gone.} No error bound survives into this deck.
\end{itemize}
\end{takeawaybox}
\end{column}
\begin{column}{0.42\textwidth}
\begin{conceptbox}[title=\textbf{Bridge to 7.B3}]\small
Two things are now owed.
\medskip
\textbf{First}, the networks obey no economics: nothing makes a value function concave or a policy monotone, though theory says both. \textbf{ICNN} and monotone networks impose it by construction.
\medskip
\textbf{Second}, with no contraction bound left, \textbf{how do you know any of this worked?}
\end{conceptbox}
\end{column}
\end{columns}
\end{frame}

\end{document}
